% 04f-bridge.tex — v3 §3.6 (second part): Prop 3.14 (BR), the X0 band (input of 04x-x0band.tex), (Osc), Lemma 3.15,
% Prop 3.16 (G1').
\subsection{The bridge, the jump at \texorpdfstring{$X_0$}{X0}, and positivity of the sharp reference}\label{sec:noise-bridge}
% v3: §3.6 (Proposition 3.14 – Proposition 3.16)

\begin{proposition}[the bridge: operator versus sharp]\label{prop:br}
% v3: Proposition 3.14
Let $o_u$ be the supremum of $\osc(\Tc)$ over windows of length $2.04H(u)$ meeting $[u,2u)$: the $H$-window oscillation of
the jump-free source. Only the kernels $C_n$ with $n\ge2$ use it, so it is read at heights $\le x/2+1.1H(x)$. For $x\ge2\Xm$
(in particular at every point of $W_x$ for $x\ge2X_0$, since $x-H(x)/2\ge2X_0-H(2X_0)/2>2\Xm$; Theorem~\ref{thm:g2} and
Proposition~\ref{prop:pointwise}(b) use it there) split
% Logged correction 90 (CORRECTIONS-vs-v3.md): stated for x >= 2X_0 but applied at the points of W_x; the proof needs only
% x >= 2X_0^- (starts >= X_0 for Corollary kr-weights(b), and x/n >= X_0^- for every level used).
\[
  (\Rres T')(x)-\rhosh(x)=\mathrm{BR}_c(x)+J_{X_0}(x),\qquad
  J_{X_0}(x):=\jz\Big[(\Rres\mathbf 1_{[X_0,\infty)})(x)-\sum_nc_n(x)\mathbf 1[x/n\ge X_0]/n\Big],
\]
where $\mathrm{BR}_c:=\Rres\Tc-\sum_nc_n\Tc(\cdot/n)/n$. Then
\[
  \log x\cdot|\mathrm{BR}_c(x)|\le(5.1C_N+C_S)\sum_{u\text{ dyadic}}\log(2u)\,o_u+8.4\,C_NE_\theta=:E_{\mathrm{BR}} .
\]
\end{proposition}

\begin{proof}
Write $\mathrm{BR}_c(x)=\sum_n\frac1n\big[\int\Tc\,dC_n^{\mathrm{nom}}-c_n\Tc(x/n)\big]+\sum_n\frac1n\int\Tc\,d\tilde E_n$, in the
notation of \S\ref{sec:kernels}.
\begin{itemize}
\item \emph{The Kantorovich--Rubinstein part.} $\int\Tc\,dC_n^{\mathrm{nom}}-c_n\Tc(x/n)=\int(\Tc-\Tc(x/n))\,dC_n^{\mathrm{nom}}$, an
  integral over $[x/n\pm R_n]$ with $R_n\le1.02H(x/n)$. Let $u$ be the dyadic scale with $x/n\in[u,2u)$; then $u\ge X_0/2$,
  since $x/n\ge\Xm$. Each half $[x/n-R_n,x/n]$, $[x/n,x/n+R_n]$ has length $\le1.02\cdot2^{\eps}H(u)\le2.04H(u)$ and meets
  $[u,2u)$, so the integral is at most $2o_u\TV(C_n^{\mathrm{nom}})$. Summing over $n$ with Corollary~\ref{cor:kr-weights}(b)
  (start $x$), this part is at most $\sum_u2o_uW^{\mathrm{nom}}_x(u)\le2C_N\sum_u\log(2u)\,o_u/\log x$, which is inside the
  first term.
\item \emph{The near-threshold part.} For $\tau$-near steps, by the Duhamel form in Proposition~\ref{prop:true-kernels} and
  $|\Tc|\le2$, it is at most $4.2\sum_d(\TV(C_d)/d)\,\theta(x/d)\le8.4C_NE_\theta/\log x$ (Lemma~\ref{lem:tau-near}). For
  $X_0$-near (terminal) steps, $\Tc$ is within $o_{X_0}$ of $\Tc(X_0)=0$ on the support, so their part is at most
  $o_{X_0}\cdot2C_N\cdot1.52\log(2X_0)/\log x$, which is inside the first term.
\end{itemize}
Only the oscillation of the smooth source enters, never the $H$-scale roughness of the noise. The jump
$\jz\mathbf 1_{[X_0,\infty)}$ has no small oscillation, and its whole effect is $J_{X_0}$, which is treated next.
\end{proof}

% 04x-x0band.tex -- v3 §3.6: the jump of the target at X_0 (Lemma 3.14A = XB, its three ranges, Lemma 3.14B = SI(Λ)).
% Input from 04-noise (bridge subsection), after Proposition 3.14.
\subsubsection{The jump of the target at $X_0$}\label{sec:x0band}
% v3: §3.6, Lemmas 3.14A and 3.14B

The source $T'$ vanishes below $X_0$ and equals $\jz=1-\rhodet(X_0)$ just above it (\S\ref{sec:noise-notation}). A jump
cannot be handled by the oscillation bounds of Proposition~\ref{prop:br}: the kernels $C_n^x$ with $x/n$ within about
$H(X_0)$ of $X_0$ straddle it, and their weight in a dyadic block is of order $\log X_0/\log x$, which is not small
compared with the margin: $\rhosh(x)\log x\ge0.97-o(1)$ for $x\ge X_0^{1.01}$, but only $\ge0.82-o(1)$ on
$[X_0+H(X_0),X_0^{1.01}]$ (Proposition~\ref{prop:g1} and Lemma~\ref{lem:margin}(d)).
% Logged correction 47: piecewise margin
We therefore treat the jump separately and exactly. Recall
$J=\jz\mathbf 1_{[X_0,\infty)}$, $\Tc=T'-J$, and put
\[
  \rhosh_1:=(1+\Kpt)^{-1}\mathbf 1_{[X_0,\infty)},\qquad\text{so that}\qquad \rhosh=\rhosh_{\mathrm c}+\jz\rhosh_1,
  \qquad J_{X_0}=\jz\bigl(\Rres\mathbf 1_{[X_0,\infty)}-\rhosh_1\bigr),
\]
with $\rhosh_{\mathrm c}:=\sum_n c_n\Tc(\cdot/n)/n$ and $J_{X_0}$ as in Proposition~\ref{prop:br}. Throughout this
subsection
\[
  \Lambda:=\log X_0\ge10^8,\qquad \delta:=H(X_0)/X_0=e^{8(\log\Lambda)^2-\Lambda},\qquad \tau_0:=\log(4.01X_0),\qquad N:=y/X_0 .
\]
By \eqref{eq:RH}, $\eps\le0.03$. Near $X_0$ the jitter width is the standard one, and its local exponent
$\eps_*:=\sup_{X_0\le v\le X_0^2}d\log H/d\log v\le16\log\Lambda/\Lambda\le2.95\cdot10^{-6}$ ($16\log\log v/\log v$ is decreasing), so that
$H(vw)/H(v)\le w^{\eps_*}$ for $X_0\le v\le vw\le X_0^2$. By Huxley's theorem \eqref{eq:huxley}, $\jz\le1.0001/\Lambda$
(\S\ref{sec:noise-notation}).

\begin{lemma}[The $X_0$ jump in exact chain form]\label{lem:xb}
% v3: Lemma 3.14A (Lemma XB)
% Logged correction 90 (CORRECTIONS-vs-v3.md): stated for y >= 2X_0, but used at every y in W_x with x >= 2X_0 (eq:jX0), so down
% to 2X_0 - H/2. The proof uses no lower bound on y; below 2X_0^- the left side of eq:xb vanishes.
Let $y\ge2\Xm$. For a level $n$ write $H_n:=H(y/n)$. For a pair $(d,w)$ with $d,w\in\NN$, $d,w\ge2$, $n=dw$, let
$D_d:=C_d^{y,\nom}-c_d\,\delta_{y/d}$ (a measure of mass $0$; $C^{y,\nom}_d$ as in \S\ref{sec:kernels-steps}), and define
\begin{align*}
  \mathbb B_0&:=\{n:\ |y/n-X_0|\le H_n/2\},\\
  \mathbb B_1&:=\{(d,w):\ \text{the step $d\to n$ is $X_0$-near, i.e. } |y/(dw)-\Xm|\le R_d/w\},\\
  \mathbb B_2&:=\{(d,w)\notin\mathbb B_1:\ |y/n-X_0|\le(1+\vartheta_{X_0})H_n/2+R_d/w\}.
  % CORRECTIONS #29 (factor 1+ϑ_{X₀})
\end{align*}
Then, on $\Good(<y)$,
\begin{equation}\label{eq:xb}
  \bigl|\Rres\mathbf 1_{[X_0,\infty)}(y)-\rhosh_1(y)\bigr|\le\frakB(y)+\frac{4.3\,C_N E_\theta}{\log y},
\end{equation}
\begin{equation}\label{eq:frakB}
  \frakB(y):=\sum_{n\in\mathbb B_0}\frac{|c_n|}{2n}+\sum_{(d,w)\in\mathbb B_1}\frac{\TV(D_d)}{2dw}
  +\lambda_0\sum_{(d,w)\in\mathbb B_2}\frac{\KRn{D_d}}{dw\cdot wH_n}.
\end{equation}
By Lemma~\ref{lem:tau-near}, $E_\theta\le10^4\Lambda^2e^{8(\log\Lambda)^2-\Lambda}$, so with $C_N\le158$
\begin{equation}\label{eq:tau-term-bound}
  4.3\,C_NE_\theta\le4.3\cdot158\cdot10^4\,\Lambda^2e^{8(\log\Lambda)^2-\Lambda}\le e^{-0.99\Lambda}\qquad(\Lambda\ge10^8),
\end{equation}
since $15.8+2\log\Lambda+8(\log\Lambda)^2\le0.01\Lambda$. If no step from the start $y$ is
$\tau$-near, the last term of \eqref{eq:xb} may be omitted.
\end{lemma}

\begin{proof}
All kernels below start at $y$; write $C_n=C_n^y$ for the true and $C_n^{\nom}$ for the nominal kernels, so that
$\Rres\mathbf 1_{[X_0,\infty)}(y)=\sum_n n^{-1}C_n([X_0,\infty))$ by \eqref{eq:R-kernels}, and
$\rhosh_1(y)=\sum_n n^{-1}c_n\mathbf 1[y/n\ge X_0]$.

\emph{Step 1: the decomposition of Proposition~\ref{prop:true-kernels}.} Write $\tilde E_n:=C_n-C_n^{\nom}$. Then
$\tilde E_n=S_n+\sum_{d<n}P^{\nom}_{d\to n}\tilde E_d$ with sources $S_n=S_n^\tau+S_n^{X_0}$ on $\tau$-near and $X_0$-near
steps, where $P_{d\to n}$, $P^{\nom}_{d\to n}$ are the step operators of \S\ref{sec:true-kernels}. By Lemma~\ref{lem:near-structure}, the $X_0$-near sources sit at terminal levels and are never
transferred, so
\[
  C_n=C_n^{\nom}+S_n^{X_0}+\tilde E_n^\tau,\qquad \tilde E_n^\tau:=\sum_{n'\mid n}\mathcal T^{\nom}_{n'\to n}S^\tau_{n'},
\]
where $\mathcal T^{\nom}_{n'\to n}$ is the nominal transfer (the signed sum over nominal chains from level $n'$ to $n$).
On a step $d\to n$ the true and nominal filters at the centre $c$ are $\varphi_\tau(c)\varphi(c)$ and
$\kappa_\tau\kappa_{X_0}$, where $\varphi(c):=\mathbf 1[c\ge\Xm]$, $\kappa_{X_0}:=\mathbf 1[y/n\ge\Xm]$, and
$\varphi_\tau,\kappa_\tau$ are the $\tau$-filters. As in the proof of Proposition~\ref{prop:true-kernels}, on $X_0$-near steps the
whole difference goes to $S^{X_0}_n$, and on the other steps the difference $(\varphi_\tau-\kappa_\tau)\kappa_{X_0}$ goes to $S^\tau_n$;
so $S_n^{X_0}=-\sum_{d\to n\ X_0\text{-near}}\mathcal Q\bigl((\varphi_\tau\varphi-\kappa_\tau\kappa_{X_0})\,D_wC_d\bigr)$, where $\mathcal Q$
places the window $\bx_{H(c)}$ at each centre $c$. Split $C_d=C_d^{\nom}+(C_d-C_d^{\nom})$ and write
$S_n^{X_0}=S_n^{X_0,\nom}+S_n^{X_0,\tau}$ accordingly.

\emph{Step 2: the nominal part, level by level.} For a centre $c$ let $U_c$ be uniform on $[-H(c)/2,H(c)/2)$ and
$G(c):=\Prob(c+U_c\ge X_0)$, so that $\int\mathbf 1_{[X_0,\infty)}\,d(\mathcal Q\mu)=\int G\,d\mu$ for every measure $\mu$ of
centres. Near $X_0$, $G(c)=\min\bigl(1,\max\bigl(0,\frac12+(c-X_0)/H(c)\bigr)\bigr)$, which is Lipschitz with constant
$\le(1+2\vartheta_{X_0})/H_n\le\lambda_0/H_n$ on the support of $D_wD_d$ (width variation, \S\ref{sec:kernels-steps}), and
$G(y/n)=\Prob(y/n+U\ge X_0)$ with $U$ uniform on $[-H_n/2,H_n/2)$. By \eqref{eq:kernel-rec} for
the nominal kernels and $D_wC_d^{\nom}=c_d\delta_{y/n}+D_wD_d$,
\[
  C_n^{\nom}([X_0,\infty))=c_n\,G(y/n)-\sum_{d\mid n,\,d<n}\kappa_{d\to n}\int G\,d(D_wD_d).
\]
Adding
\[
  S_n^{X_0,\nom}([X_0,\infty))=-\sum_{X_0\text{-near}}\int G\,(\varphi_\tau\varphi-\kappa_\tau\kappa_{X_0})\,d(D_wC_d^{\nom}),
\]
and using $(\varphi_\tau\varphi-\kappa_\tau\kappa_{X_0})\,c_d\delta_{y/n}=0$ (the filters at the nominal point equal their nominal values),
\begin{equation}\label{eq:level-defect}
\begin{split}
  &C_n^{\nom}([X_0,\infty))+S_n^{X_0,\nom}([X_0,\infty))-c_n\mathbf 1[y/n\ge X_0]\\
  &\qquad=c_n\bigl(G(y/n)-\mathbf 1[y/n\ge X_0]\bigr)-\sum_{d\mid n,\,d<n}\int G\psi_{d\to n}\,d(D_wD_d),
\end{split}
\end{equation}
where $\psi_{d\to n}:=\varphi_\tau\varphi$ on $X_0$-near steps and $\psi_{d\to n}:=\kappa_\tau\kappa_{X_0}$ (a constant) on the
others. Since $D_1=0$, only $d\ge2$ occurs.
\begin{itemize}
\item $|G(y/n)-\mathbf 1[y/n\ge X_0]|\le\frac12$, and it vanishes unless $|y/n-X_0|\le H_n/2$, i.e. unless $n\in\mathbb B_0$.
\item $D_wD_d$ has mass $0$, total variation $\TV(D_d)$, and support of radius $R_d/w$ around $y/n$. The function
  $G\psi_{d\to n}$ takes values in $[0,1]$, so $|\int G\psi_{d\to n}\,d(D_wD_d)|\le\frac12\TV(D_d)$ (subtract the constant $\frac12$).
\item If the step is not $X_0$-near, then $G\psi_{d\to n}=\kappa_\tau\kappa_{X_0}G$, and $G$ is Lipschitz with constant
  $\le\lambda_0/H_n$; by Lemma~\ref{lem:kr-facts}(K2), $|\int G\,d(D_wD_d)|\le\lambda_0\KRn{D_wD_d}/H_n=\lambda_0\KRn{D_d}/(wH_n)$. If moreover
  $|y/n-X_0|>(1+\vartheta_{X_0})H_n/2+R_d/w$, the support lies where $G$ is constant (there $H(c)\le(1+\vartheta_{X_0})H_n$)
  and the term vanishes.
  % CORRECTIONS #29
\end{itemize}
Multiplying \eqref{eq:level-defect} by $1/n=1/(dw)$ and summing over $n$ gives at most $\frakB(y)$.

\emph{Step 3: the $\tau$-part.} What remains is
\[
  \sum_n n^{-1}\bigl(\tilde E_n^\tau+S_n^{X_0,\tau}\bigr)([X_0,\infty)),
\]
which is at most $\sum_n\TV(\tilde E_n^\tau)/n+\sum_n\TV(S_n^{X_0,\tau})/n$. As in the proof of Proposition~\ref{prop:true-kernels}
(the $\infty$-norm step), the first sum is $\le2.04\lambda_0\sum_d(\TV(C_d)/d)\,\theta(y/d)\le4.1\,C_NE_\theta/\log y$, by
Corollary~\ref{cor:kr-weights}(c) at every start, the bound on $W^{\mathrm{op,nt}}_y$ in Proposition~\ref{prop:true-kernels}
(the parents of $\tau$-near steps are non-terminal), and Lemma~\ref{lem:tau-near}. The
second is $\le\sum_d(\TV(\tilde E_d)/d)\,\eta_{X_0}(y/d)$ with $\eta_{X_0}$ as in the proof of
Proposition~\ref{prop:true-kernels}; with the block bound $\sum_{y/d\in[v,2v)}\TV(\tilde E_d)/d\le2\lambda_0C_N^2E_\eta\log(2v)/\log y$
from step (ii) of that proof and $\sum_v\log(2v)\sup\eta_{X_0}\le1.52\log(2X_0)$, it is
\[
  \le3.1\,C_N^2\tau_0E_\eta/\log y\le0.1\,C_NE_\theta/\log y,
\]
because $\tilde\eta\le\theta/\log^2\Xm$ pointwise. Together,
the $\tau$-part is at most $4.3\,C_NE_\theta/\log y$; if no step is $\tau$-near, then $S^\tau=0$, $\tilde E^\tau=0$ and
$C_d=C_d^{\nom}$ for every parent of an $X_0$-near step, so the $\tau$-part vanishes.
\end{proof}

\begin{remark}
% v3: remark after Lemma 3.14A
A cruder bound by the level sum $\sum_{n}(\TV(C_n)+|c_n|)/n$ over the band does not suffice: at $n_0=6$ it is about
$0.74$, while $\frakB(6X_0)\le5/12$ (Proposition~\ref{prop:xb-ranges}). The gain comes from three facts used above: the
direct chain ($d=1$) has no centre smear; the defect of a mass-$0$ measure against a $[0,1]$-valued function is at most
half its total variation; and the window average contributes only half a jump.
\end{remark}

\paragraph{Two counting tools.} The ranges below use two elementary bounds. The first is a local count on the realised
semigroup.

\begin{lemma}[Reciprocal sums over short intervals]\label{lem:short-recip}
% v3: Lemma 3.14A, middle range (Corollary 4.2, Lemma 4.18, Corollary 1.4); fix-Bprime §12.2 counting input
On $\Good(<y)$, for every $1\le M\le y/2$ and $0<a\le0.01$,
\[
  \sum_{e\in\NN\cap M[1-a,1+a]}\frac1e\le28\,a+7M^{-3/4}.
\]
\end{lemma}

\begin{proof}
Let $I:=M[1-a,1+a]\subset[M/2,2M]$; it suffices to show $\#(\NN\cap I)\le13.5|I|+K^*(M)$ with $K^*(M)/M\le6M^{-3/4}$,
since then $\sum_{e\in I}1/e\le(27aM+K^*(M))/(M(1-a))\le28a+7M^{-3/4}$.
\begin{itemize}
\item If $M<X_0/4$, then $I\subset[1,\Xm)$, where $\NN=\N$, and $\#(\NN\cap I)\le|I|+1$; take $K^*(M):=1$.
\item If $M\ge2X_0$, cut $I\subset[P,4P]$ ($P=M/2\ge X_0$) into $\lceil|I|/6H(P)\rceil$ pieces $J$ of length $\le6H(P)$.
  % Logged correction 37: Corollary localcount(ii) at the non-dyadic scale P = M/2; the range X_0 <= M < 2X_0 moves to the third case
  By Corollary~\ref{cor:localcount}(ii), $N(J)\le T(J)+11.4|J|+K_{\mathrm{loc}}(P)+K_{\mathrm{loc}}(2P)$, and under the target \eqref{eq:target},
  $T(J)\le2|J|+2+\#\Abs(J)$. Since $K_{\mathrm{loc}}\le H/100$, summing over the pieces gives
  $\#(\NN\cap I)\le13.42|I|+K_{\mathrm{loc}}(P)+K_{\mathrm{loc}}(2P)+2+\#\Abs(I)$. By Lemma~\ref{lem:lemmaP}, each window $[u,u+h(u)]$,
  $h(u)=u^{\theta/2}$, holds at most $h(u)L_u^{-4}$ absorbed elements; covering $I$ by $\lceil|I|/h(M/2)\rceil$ such windows,
  $\#\Abs(I)\le0.01|I|+(2M)^{\theta/2}$. As $K_{\mathrm{loc}}(P)+K_{\mathrm{loc}}(2P)+2\le M^{1/8}$ and $\theta<\frac12$, take $K^*(M):=2M^{1/4}$.
\item If $X_0/4\le M<2X_0$, apply the first case to $I\cap[1,\Xm)$ and the second, with $P:=X_0$, to
  $I\cap[X_0,\infty)\subset[X_0,4X_0]$; in between,
  $\NN\cap[\Xm,X_0)$ consists of integers and of at most $2H(X_0)$ jittered primes (\S\ref{sec:noise-notation}). This gives
  $K^*(M):=4X_0^{1/4}\le4(4M)^{1/4}$.
\end{itemize}
In all cases $K^*(M)/M\le6M^{-3/4}$.
\end{proof}

The second is a Rankin bound for smooth elements, valid at every height.

\begin{lemma}[Rankin bound]\label{lem:rankin-crude}
% v3: Lemma 3.14B, proof ("Rankin gives ..."); fix-Bprime §12.3 (B)
% Logged correction 73: on Good(<z') the hypothesis holds on [1,z] only for z below the horizon of z' (Remark rem:horizon).
Assume \eqref{eq:HBS} on $[1,z]$; for instance, work on $\Good(<z')$ for some $z'$ with $z\le\check z'$
(Remark~\ref{rem:horizon}; cf.\ Lemma~\ref{lem:cm}). Let $z\ge e^{10}$, $\ell\ge0$, and $v(\ell,z):=\sum_{n\in\NN\cap(e^\ell,2e^\ell],\,\Pplus(n)\le z}1/n$. Then $v(\ell,z)\le e^{14}\,(\log z)\,e^{-\ell/\log z}$.
\end{lemma}

\begin{proof}
Put $\beta:=1/\log z$. Since $\NN$ is the free semigroup on $\PP$,
$v(\ell,z)\le e^{-\beta\ell}\sum_{\Pplus(n)\le z}n^{\beta-1}=e^{-\beta\ell}\prod_{p\le z}(1-p^{\beta-1})^{-1}$. For $p\le z$,
$p^\beta\le e$, and $e^{\beta\log p}-1\le(e-1)\beta\log p$. So
\[
  \log\prod_{p\le z}(1-p^{\beta-1})^{-1}\le\sum_{p\le z}\frac1p+(e-1)\beta\sum_{p\le z}\frac{\log p}p
  +\sum_{p}\frac{p^{2\beta-2}}{1-p^{\beta-1}}.
\]
By Lemma~\ref{lem:cm}(iv) and (iii), the first two sums are $\le\log\log z+4.1$ and $\le(e-1)(1+0.43/\log z)\le1.8$; the
last is $\le1.3$: the primes $p\le7$ contribute $\le0.95$ (using $p^\beta\le p^{1/10}$), and those with $p\ge11$
contribute $\le1.33\,e^2\sum_{p\ge11}p^{-2}\le0.31$. Hence the product is $\le e^{7.2}\log z$, and the claim follows with
room.
\end{proof}

\paragraph{The three ranges.} Put
\begin{equation}\label{eq:jX0}
  j_{X_0}(x):=\sup_{y\in W_x}\log y\cdot\jz\Bigl(\frakB(y)+\frac{4.3\,C_NE_\theta}{\log y}\Bigr),
\end{equation}
so that $\log y\cdot|J_{X_0}(y)|\le j_{X_0}(x)$ for every $y\in W_x$, by Lemma~\ref{lem:xb}.

\begin{proposition}[The $X_0$ band in three ranges]\label{prop:xb-ranges}
% v3: Lemma 3.14A, "The three ranges"
Let $\Lambda\ge10^8$, $x\ge2X_0$ and $N:=x/X_0$. On $\Good(<x)$:
\begin{enumerate}[label=(\roman*)]
\item (small) if $N\le e^{0.01\Lambda}$, then $j_{X_0}(x)\le0.4209$ (computer-assisted; see Appendix~\ref{app:cert-x0band});
\item (middle) if $e^{0.01\Lambda}\le N\le e^{600\Lambda^4}$, then $\frakB(y)\le200\delta+400N^{-1/10}$ for $y\in W_x$, and
  $j_{X_0}(x)\le e^{-0.0009\Lambda}$;
\item (huge) if $N\ge e^{600\Lambda^4}$, then $\frakB(y)\le(400\delta+e^{-60})\,\tau_0/\log N$ for $y\in W_x$, and
  $j_{X_0}(x)\le e^{-59}$.
\end{enumerate}
The three ranges cover $x\ge2X_0$ for every $\Lambda\ge10^8$.
\end{proposition}

Throughout the proof, $y\in W_x$, and we write $N$ also for $y/X_0$, which differs from $x/X_0$ by a factor
$1+O(H(x)/x)$. We use the following consequences of Theorem~\ref{thm:kr-random} for the nominal kernels (with
$H_e/H_d\le(d/e)^\eps$ from \eqref{eq:RH}, and $|c_e|\le\mu^2(e)$):
\begin{equation}\label{eq:TVD-bounds}
\begin{gathered}
  \tfrac12\TV(D_d)\le|c_d|+\tfrac{\lambda_0}4\sum_{m\mid d,\,1<m<d}\Hord(m)m^{\eps-1}|c_{d/m}|,\\
  \frac{\lambda_0\KRn{D_d}}{wH_n}\le\frac{\lambda_0^2w^{\eps-1}}4\sum_{m\mid d,\,m<d}\Hord(m)m^{\eps-1}|c_{d/m}|\quad(n=dw).
\end{gathered}
\end{equation}
For the first, $\TV(D_d)\le|c_d|+\TV(C_d^{\nom})$; for the second, $H_e/H_n\le(n/e)^\eps=(wm)^\eps$ with $m=d/e$. Since
$\lambda_0^2\le1+7\delta$, the factors $\lambda_0$, $\lambda_0^2$ change the numerical constants below by a relative amount
$\le7\delta\le e^{-9\cdot10^7}$, and we absorb them. Consequently, for every set $I$ of levels,
\begin{equation}\label{eq:TVD-sums}
\begin{gathered}
  \sum_{d\in I}\frac{\TV(D_d)}{2d}\le\sum_{d\in I}\frac{|c_d|}d+\frac{\lambda_0}4\sum_{m\ge2}\Hord(m)m^{\eps-2}\sum_{e\in I/m}\frac{|c_e|}e,\\
  \sum_{d\in I}\frac{\lambda_0\KRn{D_d}}{d\,wH_{dw}}\le\frac{\lambda_0^2w^{\eps-1}}4\sum_{m\ge1}\Hord(m)m^{\eps-2}\sum_{e\in I/m}\frac{|c_e|}e .
\end{gathered}
\end{equation}

\begin{proof}[Proof of (i): the small range]
\emph{A single level.} Let $N\le e^{0.01\Lambda}$. If $n\in\mathbb B_0$ then $|n-N|\le nH_n/(2X_0)\le0.51\,(1.01N)\delta$.
If $(d,w)\in\mathbb B_1\cup\mathbb B_2$ then $|y/n-X_0|\le H_n/2+R_d/w+(X_0-\Xm)\le1.6H(X_0)$, so $|n-N|\le2N\delta$. Since
$N\delta\le e^{0.01\Lambda+8(\log\Lambda)^2-\Lambda}<10^{-10}$, every level occurring in $\frakB(y)$ is the integer $n_0$
nearest to $N$, and $n_0<X_0$. Below $\Xm$ the semigroup is $\N$, so all divisors $d$ of $n_0$ are rational integers, all
cofactors are rigid, and no step below level $n_0$ is $\tau$-near or $X_0$-near; in particular the $\tau$-term of
\eqref{eq:xb} vanishes.

\emph{The bound $B(n_0)$.} For $2\le n$, put
\begin{gather*}
  B(n):=\frac1n\Bigl[\tfrac12\mu^2(n)+\tfrac12\sum_{d\mid n,\,1<d<n}\Bigl(2\mu^2(d)+\tfrac12\lambda\frac{X(d)}d\Bigr)\Bigr],\\
  X(d):=\sum_{e\mid d,\,1<e<d}\mu^2(e)\,e\,\Hord(d/e),\qquad\lambda:=1.0001 .
\end{gather*}
Each pair $(d,n_0/d)$ lies in at most one of $\mathbb B_1,\mathbb B_2$. By \eqref{eq:TVD-bounds} with $\eps_*$ in place of
$\eps$ and $\lambda_0(d/e)^{\eps_*}\le\lambda$ for $d\le10^6$, the $\mathbb B_1$-weight of $d$ is $\le\mu^2(d)+\frac14\lambda X(d)/d$;
the $\mathbb B_2$-weight is $\le\frac18\lambda(\mu^2(d)+X(d)/d)$, which is smaller. Hence $\frakB(y)\le B(n_0)$ for
$n_0\le10^6$.

\emph{Computation, $n\le10^6$.} The certificate of Appendix~\ref{app:cert-x0band} computes $\mu^2$, $\Hord$, $X$ and the divisor
sums exactly in integers, and evaluates $B(n)$ as a sum of at most $\tau(n)$ nonnegative terms, inflated by $1+10^{-9}$.
The maximum over $2\le n\le10^6$ is $B(6)=(\frac12+\frac12(2+2))/6=5/12$; the next values are $B(12)=0.2778$ and
$B(2)=B(4)=B(10)=\frac14$, and $\sup_{n\ge100}B(n)=0.0881$.

\emph{Analytic tail, $10^6<n\le X_0^{0.01}$.} Here we use $H_e/H_d\le(d/e)^{\eps_*}$ directly instead of $\lambda$ (all heights are $\le X_0^{1.01}$). Write
$d=em$ with $e>1$ squarefree, $m\ge2$. Then $\sum_{d\mid n}\frac1n\sum_{e\mid d}\mu^2(e)\Hord(m)m^{\eps_*-1}$ is a sum over
factorisations $n=mk$ ($k\ge2$) of $\Hord(m)m^{\eps_*-2}\cdot2^{\omega(k)}k^{-1}$, which for $\sigma:=0.2$ is at most
\[
  n^{-\sigma}\sum_{m\ge2}\Hord(m)m^{\sigma+\eps_*-2}\cdot\sup_{k\ge2}2^{\omega(k)}k^{\sigma-1}
  =n^{-0.2}\,A(1.8-\eps_*)\cdot2^{0.2},
\]
where $A(s):=1/(2-\zeta(s))-1$ and the supremum is attained at $k=2$ (since $2p^{-0.8}<1$ for $p\ge3$). In ball arithmetic,
$A(1.8-3\cdot10^{-6})\le7.4915$ (the value is $7.49142$; Appendix~\ref{app:cert-consts}). With $2^{\omega(n)}\le1.64\sqrt n$ (the maximum of $\prod_{p\mid n}2p^{-1/2}$ is at $n=6$),
\[
  B(n)\le\frac{\frac12+1.64\sqrt n}{n}+\frac14\cdot7.4915\cdot2^{0.2}\,n^{-0.2}\le\frac{\frac12+1.64\sqrt n}{n}+2.152\,n^{-0.2}
  \le0.14\qquad(n>10^6).
\]

\emph{Conclusion.} $\log y\le\Lambda+\log N+10^{-10}$, so $\log y\cdot\jz\le1.0001\,(1.01+10^{-18})\le1.010102$, and
$4.3C_NE_\theta\le e^{-0.99\Lambda}$ by \eqref{eq:tau-term-bound}. So $j_{X_0}(x)\le1.010102\cdot\frac5{12}+e^{-0.99\Lambda}\le0.42088\le0.4209$.
\end{proof}

\begin{proof}[Proof of (ii): the middle range]
Let $e^{0.01\Lambda}\le N\le e^{600\Lambda^4}$. Every set below is an interval of relative half-width $\le0.01$, so
Lemma~\ref{lem:short-recip} applies to it, and we use $|c_e|\le1$.

\emph{($\alpha$) The levels $\mathbb B_0$.} Since $H_n\le H(1.01X_0)\le1.01H(X_0)$ there, $\mathbb B_0\subset
N[1-0.51\delta,1+0.51\delta]$, so this part of $\frakB(y)$ is $\le\frac12(28\cdot0.51\delta+7N^{-3/4})\le7.2\delta+3.5N^{-3/4}$.

\emph{($\beta_1$) $X_0$-near pairs with $w\le\sqrt N$.} For fixed $w$ the condition $|y/d-w\Xm|\le R_d\le1.02H(y/d)$ and
\eqref{eq:RH} give $d\in I_w:=M_w[1-1.01\delta_w,1+1.01\delta_w]$, with $M_w:=y/(w\Xm)$ and $\delta_w:=1.03w^{\eps-1}\delta$.
By \eqref{eq:TVD-sums} and Lemma~\ref{lem:short-recip},
\[
  \sum_{d\in I_w}\frac{\TV(D_d)}{2d}\le28.3\delta_w\Bigl(1+\frac{A_0}4\Bigr)+7M_w^{-3/4}
  +\frac74M_w^{-3/4}\sum_{m\le2M_w}\Hord(m)m^{\eps-1.25}\le43\,\delta_w+41\,M_w^{-0.2},
\]
using $A_0=A_0(\eps)\le2.07$ and $\sum_{m\le X}\Hord(m)m^{\eps-1.25}\le X^{0.55}A(1.8-\eps)\le13.1X^{0.55}$, where
$A(1.77)=13.08$ (the abscissa of $\sum\Hord(m)m^{-s}$ is $1.7286<1.77$). Weighting by $1/w$ and summing over
$2\le w\le\sqrt N$, with $M_w\ge N/w$ and $\sum_{w\ge2}w^{\eps-2}\le\zeta(1.97)-1=0.674$, gives
$\le43\cdot1.03\cdot0.674\,\delta+41N^{-0.2}\cdot5N^{0.1}\le30\,\delta+205\,N^{-0.1}$.

\emph{($\beta_2$) Pairs in $\mathbb B_2$ with $w\le\sqrt N$.} Here $|y/d-wX_0|\le wH_n/2+R_d\le1.04\,w H(X_0)$, so $d$
ranges over an interval of relative half-width $\le1.05\delta$ around $y/(wX_0)$. By \eqref{eq:TVD-sums} and
Lemma~\ref{lem:short-recip}, for fixed $w$ this part is
$\le\frac14w^{\eps-1}\bigl(29.4\delta(1+A_0)+142(w/N)^{0.2}\bigr)$, and summing over $w$ gives
$\le\frac14\cdot29.4\cdot3.07\cdot0.674\,\delta+36N^{-0.2}\sum_{w\ge2}w^{-1.77}\le15.3\,\delta+34\,N^{-0.2}$.

\emph{($\gamma$) Pairs with $w>\sqrt N$.} Then $d<1.01\sqrt N$; fix $d$ and count $w$. For $X_0$-near pairs the
$w$-window $\{w:|y/d-w\Xm|\le R_d\}$ has relative half-width $\le1.02\delta(d/N)^{1-\eps}$ at height
$M':=y/(d\Xm)\ge0.98\sqrt N$, so $\sum_w1/w\le29\delta(d/N)^{1-\eps}+7M'^{-3/4}$. For pairs in $\mathbb B_2$ the window has
relative half-width $\le1.05\delta$, and by \eqref{eq:TVD-bounds} the weight carries $\frac14w^{\eps-1}\le\frac14N^{-0.485}$.
By Theorem~\ref{thm:kr-random}, $\sum_{d\le D}\TV(D_d)/(2d)\le1.52(\log D+1)$ and
$\sum_{d\le D}d^{-1}\sum_{m\mid d}\Hord(m)m^{\eps-1}|c_{d/m}|\le3.07(\log D+1)$. Hence this part is
$\le(\log N+1)\bigl(1.52(29\delta N^{-0.485}+7.2N^{-3/8})+0.77N^{-0.485}(30\delta+7.2N^{-3/8})\bigr)\le N^{-1/4}$.

\emph{Total.} $\frakB(y)\le53\,\delta+3.5N^{-3/4}+205N^{-0.1}+34N^{-0.2}+N^{-1/4}\le200\,\delta+400\,N^{-1/10}$. Finally
$\log y\le601\Lambda^4$, so $\log y\cdot\jz\le1.0001\cdot601\Lambda^3$, and with $N\ge e^{0.01\Lambda}$ and
$4.3C_NE_\theta\le e^{-0.99\Lambda}$ by \eqref{eq:tau-term-bound},
\[
  j_{X_0}(x)\le1.0001\cdot601\Lambda^3\bigl(200\,e^{8(\log\Lambda)^2-\Lambda}+400\,e^{-0.001\Lambda}\bigr)+e^{-0.99\Lambda}\le e^{-0.0009\Lambda},
\]
because $0.0001\Lambda\ge12.4+3\log\Lambda$ for $\Lambda\ge10^8$.
\end{proof}

\begin{proof}[Proof of (iii): the huge range]
Let $J:=\log N\ge600\Lambda^4$ and $W_1:=e^{100}$. The main pieces use Lemma~\ref{lem:si} below, with $J_1=300\Lambda^4$.
Recall from Lemma~\ref{lem:lemmaM} that $c_e(y)\ne0$ forces $e$ to be squarefree and $4v$-dense whenever $y/e\in[v,2v)$,
including the condition $p_1\le4v$ on its smallest prime.
% Logged correction 40: t-dense includes j = 0 (p_1 <= t)
Applying this with $v$ just above $y/(2e)$, so that $4v\le2.03\,sX_0$, an element $e$ with $y/e\le1.01\,sX_0$ is $t$-dense
with $t=4.01\,sX_0$.

\emph{($\alpha$) The levels $\mathbb B_0$.} $\mathbb B_0\subset[N(1-0.51\delta),N]\cup[N,N(1+0.51\delta)]$. Lemma~\ref{lem:si}
with $t=4.01X_0$ and $\eta=0.52\delta\in[e^{-2\Lambda},\frac12]$ gives
$\frac12\cdot2\cdot61.1(1+2.1\tau_0/J)\cdot0.52\delta\,\tau_0/(J-1)\le32\,\delta\tau_0/J$.

\emph{($\beta_1$) $X_0$-near pairs with $w\le W_1$.} As in (ii), $d\in I_w=M_w[1-1.01\delta_w,1+1.01\delta_w]$ with
$\delta_w=1.03w^{\eps-1}\delta\ge\delta e^{-100}\ge e^{-2\Lambda}$ and $\log M_w\ge J-101$. In \eqref{eq:TVD-sums}, the
element $d$ is $t_w$-dense with $t_w:=4.01wX_0$, and the element $e\in I_w/m$ is $t_{wm}$-dense with $t_{wm}:=4.01wmX_0$.
\begin{itemize}
\item For $m\le N^{1/100}/w$: the height of $I_w/m$ is $\ge N^{0.99}/1.01$, so its logarithm $J'$ is $\ge0.99J-1\ge J_1$;
  and $t_{wm}\le4.01N^{1/100}X_0\le(N^{0.99})^{1/10}$ because $X_0\le N^{1/(600\Lambda^3)}$. So Lemma~\ref{lem:si} applies, with
  $\tau_{wm}/J'\le0.0103$ and factor $1+2.1\tau_{wm}/J'\le1.022$:
  $\sum_{e\in I_w/m}|c_e|/e\le2\cdot61.1\cdot1.022\cdot1.02\,\delta_w\,\tau_{wm}/J'$, with $\tau_{wm}=\tau_0+\log w+\log m$.
\item For $m>N^{1/100}/w$: Lemma~\ref{lem:short-recip} gives $\sum_{e\in I_w/m}1/e\le29\delta_w+7$, and
  \[
    \sum_{m>N^{1/100}/w}\Hord(m)m^{\eps-2}\le(N^{1/100}/W_1)^{-0.2}A(1.77)\le e^{22.6}N^{-0.002}.
  \]
\end{itemize}
With $A_0\le2.07$ and $A_1=A_1(\eps)\le9.42$, and $J'\ge0.99J-1$,
\[
  \sum_{d\in I_w}\frac{\TV(D_d)}{2d}\le\frac{127.4\,\delta_w}{0.99J-1}\Bigl(\tau_w\bigl(1+\tfrac{A_0}4\bigr)+\tfrac{A_1}4\Bigr)
  +e^{26}N^{-0.002},\qquad \tau_w:=\tau_0+\log w .
\]
Weighting by $1/w$ and summing over $2\le w\le W_1$, with
\[
  \sum_{w\ge2}w^{-1.97}=0.674\quad\text{and}\quad\sum_{w\ge2}w^{-1.97}\log w=-\zeta'(1.97)\le1.0001
\]
(ball arithmetic; Appendix~\ref{app:cert-consts}), gives
\[
  \le\frac{132.5\,\delta}{0.99J-1}\bigl(0.674(1.518\tau_0+2.36)+1.518\cdot1.0001\bigr)+e^{31}N^{-0.002}\le138\,\delta\tau_0/J.
\]

\emph{($\beta_2$) Pairs in $\mathbb B_2$ with $w\le W_1$.} Now $d$ ranges over an interval of relative half-width
$\le1.05\delta$ around $y/(wX_0)$, and by \eqref{eq:TVD-sums} with the factor $\frac14w^{\eps-1}$ and the same use of
Lemma~\ref{lem:si} (now including $m=1$),
this part is $\le\frac14\sum_{2\le w\le W_1}w^{\eps-2}\cdot\frac{2\cdot61.1\cdot1.022\cdot1.06\,\delta}{0.99J-1}
\bigl(\tau_w(1+A_0)+A_1\bigr)+e^{31}N^{-0.002}\le72\,\delta\tau_0/J$.

\emph{($\gamma'$) Pairs with $w>W_1$.} Fix $d$; then $y/d>W_1X_0/2$. Group $d$ by the dyadic block $y/d\in[v,2v)$,
$v\ge W_1X_0/2$. For $X_0$-near pairs, the $w$-window has relative half-width $\le1.02H(v)/v\le1.02\delta(X_0/v)^{1-\eps}$
at height $\ge v/X_0$, so $\sum_w1/w\le29\delta(X_0/v)^{0.97}+7(X_0/v)^{3/4}$ by Lemma~\ref{lem:short-recip}. For pairs in
$\mathbb B_2$ the window has relative half-width $\le1.05\delta$, and the weight carries
$\frac14w^{\eps-1}\le\frac14(X_0/v)^{0.97}$. By Proposition~\ref{prop:true-kernels} and Corollary~\ref{cor:saias} (and
Corollary~\ref{cor:kr-weights} for the $m$-sums), the blocks satisfy
\begin{gather*}
  \sum_{y/d\in[v,2v)}\TV(D_d)/(2d)\le(C_N+C_S)\log(2v)/\log y\\
  \text{and}\quad\sum_{y/d\in[v,2v)}d^{-1}\sum_{m\mid d}\Hord(m)m^{\eps-1}|c_{d/m}|\le C_S\bigl((1+A_0)\log(2v)+A_1\bigr)/\log y.
\end{gather*}
Summing
the geometric series over $v=2^jW_1X_0/2$,
\begin{multline*}
  (\gamma')\le\frac{2.5\,(C_N+C_S)(\tau_0+102)}{\log y}\Bigl(29\delta\bigl(\tfrac2{W_1}\bigr)^{0.97}+7\bigl(\tfrac2{W_1}\bigr)^{3/4}\Bigr)\\
  +\frac{2.5\cdot0.77\,C_S(\tau_0+102)}{\log y}\bigl(\tfrac2{W_1}\bigr)^{0.97}(30\delta+7)\le e^{-60}\,\frac{\tau_0}J .
\end{multline*}

\emph{Total.} $\frakB(y)\le(32+138+72)\,\delta\tau_0/J+e^{-60}\tau_0/J\le(400\delta+e^{-60})\,\tau_0/\log N$. Since
$\log y\le J+\Lambda+1$ and $\tau_0\le1.0001\Lambda$, $\log y\cdot\jz\cdot\tau_0/J\le1.0003$, so with
$4.3C_NE_\theta\le e^{-0.99\Lambda}$ by \eqref{eq:tau-term-bound}, $j_{X_0}(x)\le1.0003(400\delta+e^{-60})+e^{-0.99\Lambda}\le e^{-59}$.
\end{proof}

\paragraph{Contiguity.} The three ranges meet at $N=e^{0.01\Lambda}$ and $N=e^{600\Lambda^4}$ and so cover $x\ge2X_0$ for
every $\Lambda\ge10^8$. No step uses a short-interval count at relative width below $e^{-2\Lambda}$, and no sum of the form
$\sum_w1/(w\log(N/w))$ arises: the pairs with $w>W_1$ are counted by fixing $d$.

\paragraph{A short-interval Saias bound.} The huge range rests on the following lemma. Saias' asymptotic for integers with
$t$-dense divisors has relative error $O(1/\log t)$, which is far larger than the relative widths $\eta\approx\delta$ needed
here, so it cannot be used directly. Instead we combine the joint ``$t$-dense and smooth'' profile of
Lemma~\ref{lem:bs-profile} with the short-window prime counts of Theorem~\ref{thm:lp}\eqref{eq:lpplus}.

\begin{lemma}[Short-interval Saias bound]\label{lem:si}
% v3: Lemma 3.14B (Lemma SI(Λ))
Put $\Lambda_1:=100\Lambda^3$ and $J_1:=300\Lambda^4$. Let $J:=\log N\ge J_1$, $t\in[X_0,N^{1/10}]$, $\tau:=\log t$ and
$\eta\in[e^{-2\Lambda},\frac12]$. Then, on $\Good(<y)$ for every $y\ge N$,
\begin{equation}\label{eq:si}
  S_t\bigl([N,N(1+\eta)]\bigr):=\sum_{\substack{n\in\NN\cap[N,N(1+\eta)]\\ n\ \text{squarefree, $t$-dense}}}\frac1n
  \le61.1\Bigl(1+\frac{2.1\,\tau}J\Bigr)\frac{\eta\,\tau}J .
\end{equation}
\end{lemma}

The trivial bound is about $27\eta$; \eqref{eq:si} is the Saias density $\tau/J$ inside an interval of relative width
$\eta$, with no loss in $\eta$.

\begin{proof}
Let $I:=[N,N(1+\eta)]$ and let $n\in I$ be squarefree and $t$-dense. Since $n>t$, $n$ is not prime. Put $p:=\Pplus(n)$ and
$m:=n/p$. Then $m$ is squarefree and $t$-dense with $\Pplus(m)<p\le tm$ (Theorem~\ref{thm:bs}, Step 1 of its proof); in
particular $p^2\le t\,mp=tn$, so $m=n/p\ge\sqrt{n/t}\ge\sqrt{N/t}$ (this uses the lower end $n\ge N$ of $I$). Put
$P_1:=e^{\Lambda_1}$ and split according to $p<P_1$ or $p\ge P_1$.

\emph{(B) $p<P_1$.} Then $n$ is $P_1$-smooth, and $I\subset(N/2,N]\cup(N,2N]$. By Lemma~\ref{lem:rankin-crude}, this part is
$\le2e^{14}\Lambda_1e^{-(J-1)/\Lambda_1}$. This is $\le10^{-6}\eta\tau/J$, because $\eta\ge e^{-2\Lambda}$, $\tau\ge\Lambda$, and
$(J-1)/\Lambda_1\ge3\Lambda-10^{-20}\ge2\Lambda+\log J+3\log\Lambda+34$ at $J=J_1$, while the left side grows linearly and the
right logarithmically in $J$.

\emph{(A) $p\ge P_1$: the prime sum.} For fixed $m$, $p$ runs over primes in $I/m=[X,X(1+\eta)]$ with $X:=N/m$ and
$\Lambda_p:=\log X\ge\Lambda_1-1$; in particular $X\ge X_{\mathrm{LP}}$, since $\Lambda_1-1=100\Lambda^3-1\ge10^5\Lambda^2\log^2\Lambda=\log X_{\mathrm{LP}}$.
% Logged correction 70 (CORRECTIONS-vs-v3.md): X >= X_LP made explicit (was "log X_LP = O(Lambda^2 log^2 Lambda)").
% Logged correction 45: new L_1 (log k + 3.08), the extra term Lambda_p^{-2} in Rel, the hypothesis v(1+2 eta) <= y, and the
% explicit window of Theorem lp(iii)
% Logged correction 73 (CORRECTIONS-vs-v3.md): v3 applies Theorem lp with its parameter y equal to the cutoff of Good(<y),
% where (LP-a) is not certified up to y; it is now applied with y' = 2N^{0.55}, far below the horizon.
Since $m\ge\sqrt{N/t}\ge N^{0.45}$ and $\eta\le\frac12$, also $X(1+2\eta)\le2N^{0.55}=:y'$. We apply Theorem~\ref{thm:lp}(iii)
with $y'$ in place of its parameter $y$ and $v=X$, so that its hypothesis $v(1+2\eta)\le y'$ holds; (LP-a) holds up to $y'$ on
$\Good(<y)$ by Remark~\ref{rem:horizon}, as $y'\le0.9N\le0.9y$. Theorem~\ref{thm:lp}(iii) is stated for $\Lambda_p\ge76\Lambda^3$; here
$\Lambda_p\ge\Lambda_1-1\ge76\Lambda^3$. There $\mathrm{Rel}\le0.05+\Lambda_p^{-2}\le0.1$ as soon as the exponent in
\eqref{eq:lpplus} is at least $\log(360\,G(L_1)\Lambda_p^3\eta^{-1})$. Put $\ell_p:=\log\Lambda_p$, so that
$k\le\log(1/\eta)+20\ell_p+1\le2\Lambda+20\ell_p+1$ and
$L_1=2\ell_p+\log(1/\eta)+\log k+3.08\le2\Lambda+2\ell_p+\log(2\Lambda+20\ell_p+1)+3.08$, and write
$D:=L_1+442+2\log\Lambda+\log L_1$ for the denominator of the exponent. Since $c_{\mathrm{zf}}\ge0.088/\Lambda$, the exponent
is at least $0.0792\,\Lambda_p/(\Lambda D)$.
\begin{itemize}
\item \emph{Case $\ell_p\le e^{\Lambda/2}/3$.} Then $L_1\le e^{\Lambda/2}$, so $\log G(L_1)\le\Lambda$ by Theorem~\ref{thm:lp}(iii).
  Write $\Lambda_p=76\Lambda^3r$ with $r\ge1$, and $\ell_0:=\log(76\Lambda^3)\le0.001\Lambda$. Then $\ell_p=\ell_0+\log r$; since
  $k$ and $L_1$ are at most $3\Lambda(1+\log r)$ and $1+\log r\le r$, both $\log k$ and $\log L_1$ are at most
  $\log(3\Lambda)+\log r$. So
  \[
    D\le2.003\Lambda+4\log r,\qquad \log\bigl(360\,G(L_1)\Lambda_p^3\eta^{-1}\bigr)\le3\Lambda+3\ell_p+5.9\le3.002\Lambda+3\log r,
  \]
  since $2\ell_0+446+2\log\Lambda+2\log(3\Lambda)\le0.003\Lambda$ and $3\ell_0+5.9\le0.002\Lambda$. Using $\log r\le r$ and
  $\log^2r\le r$,
  \[
  \begin{aligned}
    \Lambda D\,(3.002\Lambda+3\log r)&\le\Lambda\bigl(6.0131\Lambda^2+18.02\Lambda\log r+12\log^2r\bigr)\le6.0132\,\Lambda^3r\\
    &<0.0792\cdot76\,\Lambda^3r=0.0792\,\Lambda_p ,
  \end{aligned}
  \]
  which is the required inequality.
\item \emph{Case $\ell_p>e^{\Lambda/2}/3$.} Then $\ell_p\ge\Lambda^2$, so $k\le21\ell_p$ and
  $L_1\le2\ell_p+2\Lambda+\log(21\ell_p)+3.08$. Since $2\Lambda+457+2\log\Lambda+2\log(21\ell_p)\le\ell_p$ (with huge room, as
  $\ell_p\ge\Lambda^2\ge10^{16}$), both $L_1+450+2\log\Lambda\le3\ell_p$ and $D=L_1+442+2\log\Lambda+\log L_1\le3\ell_p$. So
  $G(L_1)\le400\cdot3\ell_p\le\Lambda_p$, hence $\log(360\,G\Lambda_p^3\eta^{-1})\le4\ell_p+2\Lambda+5.9\le5\ell_p$, and the
  required inequality $0.0792\,\Lambda_p\ge\Lambda\cdot3\ell_p\cdot5\ell_p$ holds because
  $190\Lambda\ell_p^2\le\ell_p^3\le e^{\ell_p}=\Lambda_p$.
\end{itemize}
So $\mathrm{Rel}\le0.1$ for every $\Lambda_p\ge\Lambda_1-1$. The window $\varphi$ of Theorem~\ref{thm:lp}(iii) majorises
$\mathbf 1_{[X,X(1+\eta)]}$, takes values in $[0,1]$ and is supported in $[X(1-w),X(1+\eta+w)]$ with $w=\eta\Lambda_p^{-2}$. As
$d\Pi_0\le du/\log u$, this gives
$\Pi([X,X(1+\eta)])\le\int\varphi\,d\Pi\le1.1\int\varphi\,d\Pi_0\le1.1\,(1+3\Lambda_p^{-2})\,\eta X/\log X$: the support has
length $\eta(1+2\Lambda_p^{-2})X$, and $\log u\ge\log X-2w$ on it. As
$\Lambda_p\ge10^{25}$, $1.1(1+3\Lambda_p^{-2})\le1.1001$. Hence
\[
  \sum_{p\in I/m}\frac1p\le\frac1X\,\Pi\bigl([X,X(1+\eta)]\bigr)\le1.1001\,\frac{\eta}{\log X}=\frac{1.1001\,\eta}{\log(N/m)} .
\]

\emph{(A) $p\ge P_1$: the $m$-sum.} Group $m$ into dyadic blocks $(M_k,2M_k]$, $M_k:=2^kM_0$ with $M_0:=\frac12\sqrt{N/t}$,
up to $M_k\le N(1+\eta)/P_1$. In a block, $\Pplus(m)<p\le N(1+\eta)/m<2N/M_k$ and $\log(N/m)\ge\sigma_k:=\log(N/2M_k)$, so
\[
  S_A\le1.1001\,\eta\sum_k\frac{W_t(M_k;\,2N/M_k)}{\sigma_k},\qquad
  W_t(M;z):=\sum_{\substack{m\in\NN\cap(M,2M]\\ m\ \text{sqf, $t$-dense},\ \Pplus(m)<z}}\frac1m .
\]
Put $\ell_k:=\log M_k$ and $u_k:=\ell_k/\log(2N/M_k)$. Since $\tau\le J/10$, $\ell_k\ge(J-\tau)/2-2\log2\ge\tau$. Recall from
Lemma~\ref{lem:bs-profile} the profile $U(\ell)=\ell^{1/4}\bigl(1+\sqrt{2\log(\ell/\tau)+4}\bigr)$,
$Q(\ell,u)=A+a$ with $A=11(1+u)^3\log(2+u)/\ell$ and $a=(u-\ell^{1/4})_+^2/(2\sqrt\ell)$, and $\Phi_\infty\le38.3$.

\emph{Region 1: $u_k\le U(\ell_k)$.} By Lemma~\ref{lem:bs-profile},
$W_t(M_k;2N/M_k)\le(\tau/\ell_k)\,\Phi_\infty\,\rho(u_k-1)\,e^{Q(\ell_k,u_k)}$, with $\rho(v):=1$ for $v\le1$.
\begin{itemize}
\item For $u\le\ell^{1/4}$: $a=0$ and $A\le11\cdot1.01\,\ell^{-1/4}(\frac14\log\ell+0.01)\le7\cdot10^{-7}$, because
  $\ell\ge0.449J\ge10^{34}$ (as $J\ge300\Lambda^4\ge3\cdot10^{34}$) and $\ell^{-1/4}\log\ell$ is decreasing. For $\ell^{1/4}<u\le U(\ell)$: $e^{Q}\le e^{2}\ell/\tau$ and
  $\rho(u-1)\le1/\Gamma(\ell^{1/4})$, so these $k$ contribute less than $10^{-6}\tau/J$ in total.
\item With $\tilde J:=J+\log2$ we have $u_k=\ell_k/(\tilde J-\ell_k)$ and $\sigma_k=\tilde J-\ell_k-2\log2$. The function
  $f(\ell):=\rho(u(\ell)-1)/(\ell(\tilde J-\ell))$ satisfies
  \[
    \int f(\ell)\,d\ell=\frac1{\tilde J}\int\rho(u-1)\frac{du}u\qquad(\ell=\tilde Ju/(1+u)).
  \]
  The sum over $\ell_k\in\ell_0+(\log2)\Z$ differs from $(\log2)^{-1}\int f$ by at most the total variation of $f$ on
  Region~1, which is $\le4\sup f\le10^{-6}/J$, since there $\tilde J-\ell\ge\ell/U(\ell)$.
\item The lower limit is $u_{\min}\ge(J-\tau-6)/(J+\tau+8)$, and by $\int_a^\infty\rho(w-1)\,dw/w=\rho(a)$ ($a\ge1$),
  \[
    \int_{u_{\min}}^\infty\rho(u-1)\frac{du}u=\log\frac1{u_{\min}}+\rho(1)\le1+\log\frac{J+\tau+8}{J-\tau-6}\le(1+10^{-6})\Bigl(1+\frac{2.1\tau}J\Bigr),
  \]
  using $\log\frac{1+x}{1-x}=2x(1+\frac{x^2}3+\frac{x^4}5+\cdots)\le2.1x$ for $0\le x\le0.1$.
\end{itemize}
Hence Region~1 contributes at most
$1.1001\cdot\frac{38.3}{\log2}\,(1+3\cdot10^{-6})\bigl(1+\frac{2.1\tau}J\bigr)\frac{\eta\tau}J
\le60.81\bigl(1+\frac{2.1\tau}J\bigr)\frac{\eta\tau}J$.

\emph{Region 2: $u_k>U(\ell_k)$.} Drop the density condition: by Lemma~\ref{lem:rankin-crude},
$W_t(M_k;2N/M_k)\le e^{14}\log(2N/M_k)\,e^{-u_k}$, and $\log(2N/M_k)\le1.01\sigma_k$. Since $U(\ell)\ge3\ell^{1/4}\ge3(0.449J)^{1/4}$
and there are at most $J/\log2+1$ blocks, Region~2 contributes $\le1.1001\,\eta\,e^{14}\cdot1.01\,(J/\log2+1)\,e^{-3(0.449J)^{1/4}}
\le10^{-6}\eta\tau/J$.

Adding (B), Region~1 and Region~2 gives \eqref{eq:si}, since $60.81+2\cdot10^{-6}\le61.1$.
\end{proof}

\begin{remark}
% v3: §3.6, after Lemma 3.14B
No cancellation of the M\"obius function and no short-interval asymptotic for smooth numbers is used. The smooth part
needs no density saving: its Rankin factor $e^{-J/\Lambda_1}$ beats $\eta/J$ outright. The thresholds $\Lambda_1$ and $J_1$
depend only polynomially on $\Lambda$; if the zero-free constant of Theorem~\ref{thm:lp} were only $c_{\mathrm{zf}}\ge\Lambda^{-A}$,
it would suffice to take $\Lambda_1=100\Lambda^{2+A}$ and $J_1=300\Lambda^{3+A}$, and the middle range of
Proposition~\ref{prop:xb-ranges} would still close.
\end{remark}


\paragraph{Oscillation of the target.}\label{par:osc} All bounds are for $\Tc$, and none uses a local count at the height where
the oscillation is read.
% v3: (Osc)
\begin{equation}\label{eq:osc}
  \sum_u\log(2u)\,o_u\le e^{-83},\qquad\sum_v(\log3v)^2\,o_v\le e^{-41},\qquad\sup_v\log v\cdot o_v\le e^{-120},
  \qquad E_{\mathrm{BR}}\le e^{-76}.
\end{equation}
\begin{proof}
\begin{itemize}
\item Near $u=X_0$, $\Tc$ is continuous: it is $0$ below $X_0$ and $1-\rhodet-\jz$ above, which vanishes at $X_0$.
\item $o_{\mathrm{det}}\le(2.04H(u)+2)X_0^{5/12}/u$, from the linear window of length $\kappa u$ in \eqref{eq:rhodet} and the
  unit spacing of the rigid set.
\item On $\Good(<u)$, the $\Abs$-part of $\Tc$ is $\sigma_{\Abs}=K_{\Abs}\rhot-\rho_{\Abs}$ (as $\bar\rho_{\Abs}=\rho_{\Abs}$,
  Remark~\ref{rem:stopped}). So $o_{\Abs}\le2\sup|\sigma_{\Abs}|\le4L_u^{-4}$, by Lemma~\ref{lem:rhoA}(c) (from parts (a), (b)
  and Lemma~\ref{lem:lemmaP}). Lemma~\ref{lem:lemmaP} holds on $\Good(<t)$ at heights up to $t+1.01H(t)$, using only
  $E_{\mathrm{cl}}$ at lower heights and Chebyshev below $2y$. So reading $o_{\Abs}$ at the current height $t$ costs no event
  later than $t$. Moreover $o_{\Abs}$ vanishes unless
  $\log\log u>2\log\Lambda+5.97$ ($=42.8$ at $\Lambda=10^8$), since $\Abs$ is empty below that height.
\item $o_R\le\sum_{k:\,x_k\le u}\gamma_k\kappa_kH(u)u^{-1}(x_k/u)^{1-\theta}\le H(u)u^{-\theta}$, since
  $\gamma_k\kappa_k\le x_k^{1-\theta}$ and $u\ge x_1\ge\exp(X_0^2)$.
\item Hence $\sum_u\log(2u)\,o_u\le4.2H(X_0)X_0^{-7/12}\log X_0+5e^{-2\cdot42.8}\le e^{-83}$, and
  $E_{\mathrm{BR}}\le883e^{-83}+8.4C_NE_\theta\le e^{-76}$. Also $\sum_v(\log3v)^2o_v\le5.8e^{-42.8}\le e^{-41}$ and
  $\sup_v\log v\cdot o_v\le4e^{-3\cdot42.8}\le e^{-120}$. All of these improve with $\Lambda$.
\end{itemize}
No triangular windows and no Lipschitz bound are needed.
\end{proof}

\begin{lemma}[the noise terms, in predictable form]\label{lem:predictable}
% v3: Lemma 3.15
With \eqref{eq:decomp},
\begin{equation}\label{eq:predictable}
  M(x)-(B\Rres T')(x)=(\vartheta/H)(x)-(B\Kpt B\Rres F_r)(x)-(B\Rres G)(x),
\end{equation}
and, on $\Good(<x)$,
\begin{multline*}
  \log x\cdot|M-B\Rres T'|(x)\le\eps_N:=5.1C_N\sum_u\log(2u)\bar F_r(u)+15.3C_NC_S'\sum_v(\log3v)^2F_r^+(v)\\
  +\sup_x\frac{6\log x}{H(x)}\le10^{-13},
\end{multline*}
where $\bar F_r(v):=o_v+6/H(v)+\bar\nu(v)+54\eta(v)$ bounds $|F_r|$ on $[v,2v)$, and $F_r^+:=2\bar F_r$ bounds $|(1-B)F_c|$.
Every sum is over dyadic scales $\ge X_0/2$; there is no dependence on $x$, and no event at the current scale is used.
\end{lemma}

\begin{proof}
\emph{The identity.} By \eqref{eq:pushthrough} for $F_c$ and $\Kpt B\Rres=1-\Rres$, $\Kpt\rhot_c=F_c-\Rres F_c+\Rres G$, where
$\rhot_c:=\rhot-(1+B\Kpt)^{-1}BJ$. So $M=B\Rres J+B\Tc+\vartheta/H-BF_c+B\Rres F_c-B\Rres G$. Insert $F_c=\Tc+F_r$ and
$\Rres\Tc+\Rres J=\Rres T'$. In particular, the current-scale terms $-B(\nu+\NT)(x)$ of $BF_c$ cancel exactly against those of
$B\Rres F_c$ (its $n=1$ term). Every term on the right of \eqref{eq:predictable} involves $\nu$, $\NT$ and $\rhot$ only at scales
$\le x/2+2H$.

\emph{The bounds.}
\begin{itemize}
\item $\log x\cdot|(B\Kpt B\Rres F_r)(x)|\le5.1C_N\sum_{u\le x/2+2H}\log(2u)\bar F_r(u)$. The kernels of $\Kpt B\Rres=1-\Rres$
  are the $C_n$ with $n\ge2$, so Proposition~\ref{prop:true-kernels} applies ($5.1C_N$ bounds both kinds of blocks).
\item $\log x\cdot|(B\Rres G)(x)|\le15.3C_NC_S'\sum_v(\log3v)^2F_r^+(v)$, since
  $|G(y)|\le3\sum_v\omega_y(v)\sup_{[v,3v]}F_r^+$ (Lemma~\ref{lem:S}(iii)).
\item $|\vartheta/H|\le6/H(x)$, and $\sup\rhot\le1.01$ surely (\S\ref{sec:noise-perscale}), so $|\NT|\le2.02\eta\le54\eta$ by
  \eqref{eq:NT}.
\end{itemize}
By \eqref{eq:sw-sums} (note that $\sum_u\log(2u)\bar\nu(u)$ and $\sum_v(\log 3v)^2\bar\nu(v)$ are dominated by the sums
there), Lemma~\ref{lem:tau-near} ($E_\eta^0$ and $E_\eta^{(2)}$) and \eqref{eq:osc}, $\eps_N\le10^{-13}$. The largest part is
$30.6C_NC_S'e^{-41}$, from the oscillation of $\sigma_{\Abs}$.
\end{proof}

\paragraph{Design facts used} (\S\ref{sec:construction}).
\begin{itemize}
\item $R=\sum_kR_k$, with $R_k\equiv0$ on $[1,x_k+2H(x_k)]$, $x_1\ge\exp(X_0^2)$, $x_{k+1}\ge\exp(x_k)$, and
  $\gamma_k\in[\Gamma_k,2\Gamma_k]$ with $\Gamma_k=\kappa_0x_k^{1-\theta}/(2\log^3x_k)$; so $\log\gamma_k\asymp(1-\theta)\log x_k$.
\item \eqref{eq:Rkprime}: $|R_k'(v)|\le\kappa_k(v/x_k)^{\theta-1}$ with $\kappa_k=2\kappa_k'/\log x_k$, $\kappa_k'=\kappa_0/\log^2x_k$,
  $\kappa_0\le\frac14$.
\end{itemize}

\paragraph{A squarefree average.} $F_\theta:=\sup_yy^{\theta-1}\sum_{n\in\NN,\,n\le y}\mu^2(n)n^{-\theta}\le1.33$ for $\theta\le\frac12$
(in particular for $\theta\in[\frac14,\theta_1]$):
\begin{itemize}
\item $F_\theta\le F_{1/2}$, since $y^{\theta-1}n^{-\theta}=y^{-1}(y/n)^\theta$ increases with $\theta$;
% Logged correction 54: 1.32964 (an upper bound) replaces 1.3296; the maximum is 1.3296349 (independent check R8, C2).
\item a direct computation of $F_{1/2}$ on $\N$ for $y\le10^6$ gives the maximum $1.32964$, at $y=7$ (see
  Appendix~\ref{app:cert}, \texttt{finite-n});
\item for $y\ge10^6$, the elementary bound $|Q(u)-6u/\pi^2|\le3\sqrt u$ for the squarefree count $Q$ gives
  $y^{\theta-1}\sum_{n\le y}\mu^2(n)n^{-\theta}\le(6/\pi^2)/(1-\theta)+3y^{-1/2}+1.5y^{-1/2}\log y\le1.25$; and
  \textup{(H$_{\mathrm{BS}}$)} changes the count for $\NN$ by a relative $O(\eta_1)$.
\end{itemize}
% Logged correction 9 (CORRECTIONS-vs-v3.md): explicit bound |Q(u) - 6u/pi^2| <= 3 sqrt u replaces v3's (1+o(1)).

\begin{proposition}[positivity of the sharp reference]\label{prop:g1}
% v3: Proposition 3.16
On $\Good(<x)$,
% Logged correction 47 (CORRECTIONS-vs-v3.md): v3 states 0.97 - o(1) for all x >= X_0 + H(X_0); the edge term rho_edge of
% Lemma margin(d) lowers it to 0.82 on [X_0 + H, X_0^{1.01}].
\[
  \rhosh(x)\cdot\log x\ge
  \begin{cases}
    0.82-o(1), & X_0+H(X_0)\le x\le X_0^{1.01},\\
    0.97-o(1), & x\ge X_0^{1.01}/2.
  \end{cases}
\]
\end{proposition}

The cutoff $\Xm$ in $\Kpt$ does not affect $\rhosh$: $c_n(y)=\mu(n)\chi^t_y(n)\mathbf 1[y/n\ge\Xm]$ (Lemma~\ref{lem:lemmaM}),
and $T'=0$ on $[\Xm,X_0)$. The proof has two ranges. Put $\log X_{\mathrm{sw}}:=X_0^{1/100}$. For every $\Lambda\ge10^8$,
$X_0^{1.01}\le X_{\mathrm{sw}}$, $\log X_{\mathrm{sw}}\le(0.0048/\kappa)^{1/5}$ (the limit for $\mu_0$ below, which is
$\ge X_0^{0.079}$), and $\log X_{\mathrm{sw}}=e^{\Lambda/100}\ge10^5\Lambda^2\log^2\Lambda=\log X_{\mathrm{LP}}$.
% Logged correction 69 (CORRECTIONS-vs-v3.md): X_sw >= X_LP made explicit (was ">>").

\begin{proof}[Proof for range \textup{(A)}: $X_0+H\le x\le X_{\mathrm{sw}}$]
No designed oscillation is active, since $x_1\ge\exp(X_0^2)>X_{\mathrm{sw}}$. So $T'=1-\rhodet+\sigma_{\Abs}$ there.
Lemma~\ref{lem:margin} (\eqref{eq:ratid}, \eqref{eq:shdecomp}) gives
$(1+\Kpt)^{-1}(1-\rhodet)=\mathrm{MA}[\mathbb P]+\rho_{\mathrm{edge}}+(1+\Kpt)^{-1}(\mu_0+\mu_1)$, where:
% Logged correction 47: support convention for MA[P] (primes >= X_0 only), the edge term moved out of mu_0 into its own
% item (d), and the conclusion stated piecewise.
\begingroup\emergencystretch=3em
\begin{enumerate}[label=(\alph*)]
\item $\mathrm{MA}[\mathbb P](q):=\mathrm{MA}_{\kappa q}[\text{rational primes}\ge X_0](q)$ is the density of the new primes (the
  rational primes $\ge X_0$) in the linear window at $q$ (\S\ref{sec:baserange}); it vanishes for $(1+\kappa)q\le X_0$. By
  \eqref{eq:huxley} and Lemma~\ref{lem:margin}(a) (Huxley and Heath-Brown, $X_0\ge X_{\mathrm{Hux}}$),
  $\mathrm{MA}[\mathbb P](x)\log x\ge0.99$ for every $x\ge X_0$;
\item $|\mu_0(u)|\le2\kappa(\log u)^3$ for every $u\ge X_0$, in particular on range (A)
  (Lemma~\ref{lem:margin}(b)). This smooth part contains the drift $+(\kappa/2)\Psi$, the partial windows, the decision-order
  mismatch (bounded by a direct count, $\le2.02(H(X_0)+1)(\log u+1)X_0^{-1}\sup|f_0|$; ``the rest'' in the proof of
  Lemma~\ref{lem:margin}(b)), and the $O(1/(\kappa x))$ term of \eqref{eq:ratid};
% Logged correction 70: the mismatch is described as in correction 53 (was "the order mismatch H(q)/(kappa q)"; pass 20, N40).
\item $\sum_{u\text{ dyadic}\ge X_0}\log(2u)\sup_{[u,2u]}|\mu_1|\le7\cdot10^{-6}$ \eqref{eq:mu1sum};
\item $\rho_{\mathrm{edge}}:=(1+\Kpt)^{-1}\mu_{\mathrm{edge}}$ is the edge at $X_0$, where $\mu_{\mathrm{edge}}$ is the part of
  $K_{\mathbb Q}\mathrm{MA}[\mathbb P]$ that lies below $X_0$ \eqref{eq:shdecomp}. By Lemma~\ref{lem:margin}(d),
  $\rho_{\mathrm{edge}}(x)\log x\ge-0.167$ for $X_0\le x\le X_0^{1.01}$, and
  $|\rho_{\mathrm{edge}}(x)|\log x\le3X_0^{-0.01}+X_0^{-0.34}$ for $X_0^{1.01}/2\le x\le X_{\mathrm{sw}}$.
\end{enumerate}
\endgroup
Now bound the responses.
\begin{itemize}
\item $\mu_0$: by Lemma~\ref{lem:lemmaM} ($|c_n|\le1$ and $\sum|c_n|/n\le\log x+B_1$, with $B_1=1$),
  \[
    \log x\cdot|(1+\Kpt)^{-1}\mu_0(x)|\le\log x\,(\log x+1)\sup_{u\le x}|\mu_0(u)|\le2.1\,\kappa(\log x)^5\le X_0^{-0.34}
  \]
  for $\log x\le X_0^{1/100}$, since $\kappa=X_0^{-5/12+2\eps_\kappa}=X_0^{-0.39666\ldots}\le X_0^{-0.3966}$.
% Logged correction 69 (R13, C2): "kappa <= X_0^{-0.3967}" was false as written; 2.1 X_0^{-0.3966} X_0^{0.05} <= X_0^{-0.34} still.
\item $\mu_1$: by Corollary~\ref{cor:saias}, $\log x\cdot|(1+\Kpt)^{-1}\mu_1(x)|\le C_S\cdot7\cdot10^{-6}\le6\cdot10^{-4}$.
\item $\sigma_{\Abs}$: likewise, at most $C_S\sum_v\log(2v)\cdot2L_v^{-4}\mathbf 1[\log\log v>2\log\Lambda+5.97]\le e^{-80}$.
\end{itemize}
So, pointwise, for $X_0+H\le x\le X_0^{1.01}$,
\[
  \rhosh(x)\log x\ge0.99-0.167-6\cdot10^{-4}-X_0^{-0.34}-e^{-80}\ge0.822,
\]
and for $X_0^{1.01}/2\le x\le X_{\mathrm{sw}}$,
\[
  \rhosh(x)\log x\ge0.99-6\cdot10^{-4}-2X_0^{-0.34}-3X_0^{-0.01}-e^{-80}\ge0.989\ge0.97,
\]
for every $X_0\ge\max(e^{10^8},X_{\mathrm{Hux}})$; no constant here depends on $\Lambda$. By Lemma~\ref{lem:margin}(d)(i), the loss
$0.167$ occurs only on the windows $N\cdot(X_0/(1+\kappa),X_0)$ with $\mu(N)=1$, each of relative length $\kappa$.
\end{proof}

\begin{proof}[Proof for range \textup{(B)}: $x\ge X_{\mathrm{sw}}$]
Call $\gamma_j$ \emph{resolvable at $x$} if $\gamma_j\le\Gamma(x):=e^{c_{\mathrm{LP}}\sqrt{\log x}/4}$, and let $\gamma_*$ be the largest
resolvable $\gamma_j$ ($\gamma_*:=1$ if there is none). Put
\[
  \delta:=\min\big(\gamma_*^{-2}\log^{-4}x,\ X_0^{-1/2}\big),\qquad J:=[x(1-2\delta),\,x(1-\delta)].
\]
% Logged correction 69: the two comparisons below made explicit (were "delta x >> H" and "c_LP sqrt(x_j) >> log x_j").
The window $J$ ends $\delta x\ge x^{0.98}>H(x)$ before $x$ (as $1/\delta\le x^{0.02}$, Step~3), so the proof uses only the realised
data before $x-H$, on $\Good(<x)$. In particular Theorem~\ref{thm:lp} is applied with $y:=x(1-\delta)\le x-H(x)$, below the
certified horizon of $\Good(<x)$, where (LP-a) holds (Remark~\ref{rem:horizon}).
% Logged correction 73: the local prime number theorem is applied below the certified horizon of Good(<x).
% Logged correction 84 (R21, 84-1): k is defined here, before its first use, and only for x >= x_1.
For $x\ge x_1$ let $k:=\max\{j:x_j\le x\}\ge1$; for $x<x_1$ no $R_j$ is active on $[1,x]$ (Step~4). All old oscillations $j<k$ are resolvable: $\log\gamma_j\le(1-\theta)\log x_j\le\log x_j$ (as
$\gamma_j\le2\Gamma_j\le x_j^{1-\theta}$), while $\log x\ge x_{k-1}\ge x_j$ and $x_j\ge\exp(X_0^2)$, so
$c_{\mathrm{LP}}\sqrt{\log x}/4\ge(0.0044/\Lambda)\sqrt{x_j}\ge\log x_j$.

\emph{Step 1: the window average of $\rhot$.}
\begin{itemize}
\item $\int_J\rhot=\#\{q:s_q\in J\}\pm(H(x)+2)$, the error coming from the partial boxes at the two ends of $J$; and slots lie
  within $H/2$ of values.
% Logged correction 31 (CORRECTIONS-vs-v3.md): the same partial-box effect; v3 has +-2. Absorbed by the +-27H(x) below.
  By the local count, and since the prime powers
  $p^k\in J$ with $k\ge2$ number at most $2\sqrt x$, $\int_J\rhot=\Pi(J)\pm(27H(x)+2K_{\mathrm{loc}}+6+2\sqrt x)$.
\item By Theorem~\ref{thm:lp}(ii), $\Pi(J)=\delta x\,\pi_0(x)(1+\eps_J)+O(Kxe^{-c_{\mathrm{LP}}\sqrt{\log x}})$ with
  $0\le\eps_J\le2.01\,\delta/\log x$, where $\pi_0(x):=(1-1/x)/\log x$ (which decreases, and $J\subset[x(1-2\delta),x]$).
\item Since $1/\delta\le\Gamma^2\log^4x+X_0^{1/2}$ with $\Gamma^2\le e^{c_{\mathrm{LP}}\sqrt{\log x}/2}$, the error term is at most
  \[
    \pi_0(x)\bigl(2K\log^5x\,e^{-c_{\mathrm{LP}}\sqrt{\log x}/2}+2K\log x\,X_0^{1/2}e^{-c_{\mathrm{LP}}\sqrt{\log x}}\bigr);
  \]
  and
  $c_{\mathrm{LP}}\sqrt{\log x}\ge2(\log K+\Lambda+10\log\sqrt{\log x})+10$, as $\log K\le98+5.9\log\Lambda\le\Lambda$ and
  $(0.0176/\Lambda)\sqrt{\log x}\ge4\Lambda+20\log\sqrt{\log x}+10$ for $\sqrt{\log x}\ge X_0^{0.005}$. So the error is at most
  $\pi_0(x)\,e^{-\Lambda}$, and $\mathrm{avg}_J\rhot=\pi_0(x)+o(1/\log x)$.
% Logged correction 70 (check R13, C4): the condition is stated against both error terms (c_LP sqrt(log x)/2 must beat
% log(2K) + 5 log log x); the earlier form covered only log log x <= 0.6 Lambda.
% Logged correction 70 (CORRECTIONS-vs-v3.md): the comparison with log K made explicit (was "recall log K = O(Lambda)").
\end{itemize}

\emph{Step 2: $\rhosh$ against $\rhot$ on average.}
\begin{itemize}
\item $\rhot-\rhosh=[\rhot+\nu+\NT-B\Rres T']-\nu-\NT+B[\Rres T'-\rhosh]+(B\rhosh-\rhosh)$.
\item Lemma~\ref{lem:predictable} and Proposition~\ref{prop:br} bound the first and fourth terms by
  $(\eps_N+E_{\mathrm{BR}}+j_{X_0}(x))/\log x$. In range (B), $\log(x/X_0)\ge X_0^{1/100}-\Lambda\ge600\Lambda^4$, so $x$ lies in
  the huge range of Proposition~\ref{prop:xb-ranges}, where $j_{X_0}\le e^{-59}$. Also $|\nu|\le\bar\nu$ and $|\NT|\le54\eta$.
\item Averaged over $J$, $\int_J(Bf-f)$ is a boundary term $\le2H(x)\sup|f|$, and $\sup|\rhosh|\le2(\log x+B_1)$.
\item Hence $\mathrm{avg}_J\rhosh\ge\pi_0(x)-(\eps_N+E_{\mathrm{BR}}+e^{-59}+o(1))/\log x$.
\end{itemize}

\emph{Step 3: pointwise against average, for all components except the current oscillation.} For $y\in J$,
$\rhosh(x)-\rhosh(y)=\sum_n[c_n(x)T'(x/n)-c_n(y)T'(y/n)]/n$.
\begin{itemize}
\item \emph{Threshold changes:} at most $2\cdot$(Lemma~\ref{lem:threshold-mass} at $2\delta$) $=o(1/\log x)$, since $|T'|\le2$,
  $\delta\le\log^{-4}x$ and $\log H(x)=8(\log\log x)^2$.
% Logged correction 59: the hypothesis 6H(x)x^{-1/3} <= 2 delta <= 1/4 of Lemma threshold-mass, checked here.
  The lemma applies at $2\delta$: $2\delta\le2X_0^{-1/2}\le\frac14$; and $1/\delta\le\Gamma(x)^2\log^4x+X_0^{1/2}\le x^{0.02}$ (as
  $\gamma_*\le\Gamma(x)$ and $\log x\ge X_0^{1/100}$), while $6H(x)x^{-1/3}\le x^{-0.3}$ by Theorem~\ref{thm:caplw}(c), so
  $2\delta\ge6H(x)x^{-1/3}$.
\item \emph{The jump of $T'$ at $X_0$:} only $n\in[y/X_0,x/X_0]$ see it. Put $N_*:=x/X_0$, $I:=[N_*(1-2\delta),N_*]$ and
  $L_I:=\log(N_*/2)$. By the lag bound,
% Logged correction 66: renamed from L_*, which clashes with L_* = 10^7 of eq:absorbed.
$|N-T|\le H(N_*)$ at both ends of $I$, and under the target \eqref{eq:target},
  $T(I)\le2.0001|I|+1+\#\Abs(I)$ (as in Corollary~\ref{cor:localcount}(i)). Cover $I$ by $\lceil|I|/h(N_*(1-2\delta))\rceil$ consecutive windows $[u,u+h(u)]$,
  $h(u)=u^{\theta/2}$; by Lemma~\ref{lem:lemmaP} each holds at most $h(u)L_u^{-4}$ absorbed elements, so
  $\#\Abs(I)\le(2.01\,\delta N_*+N_*^{\theta/2})L_I^{-4}$. Hence $\#(\NN\cap I)\le4.01\,\delta N_*+2H(N_*)+2+N_*^{\theta/2}$, and
  this part is at most $\jz\sum_{n\in\NN\cap I}1/n\le\jz\big(4.1\,\delta+3H(N_*)/N_*+2N_*^{\theta/2-1}\big)=o(1/\log x)$.
% Logged correction 14 (CORRECTIONS-vs-v3.md): v3 omits #A(I) here (same omission as correction 4); with T(I) <= 2|I|+2+#A(I)
% the constant 2.1 delta becomes 4.1 delta (noted under #14).
% Logged correction 93 (CORRECTIONS-vs-v3.md): for y in J the level points x/n and y/n differ by up to 2 delta x/n, not
% delta x/n; the smooth-change bounds for rho_det and for the resolvable R'_j are restated at displacement 2 delta.
\item \emph{$\rhodet$, smooth change:} for $y\in J$ and $x/n\ge X_0$ put $v:=x/n$ and $v':=\max(y/n,X_0)$, so that $|v-v'|\le2\delta v$
  (if $y/n<X_0$, then $T'(x/n)-T'(y/n)=\jz+\rhodet(X_0)-\rhodet(x/n)$, as $R'$ and $\sigma_{\Abs}$ vanish near $X_0$, and $\jz$ is
  the previous item); then
  $|\rhodet(v)-\rhodet(v')|\le\min\big(2(2\delta(1+\kappa)v+1)/(\kappa v),\ 3\Psi(2v,X_0)/(\kappa v)\big)$. Indeed the windows
  $[v,(1+\kappa)v]$ and $[v',(1+\kappa)v']$ overlap (as $2\delta<\kappa$) and differ by two intervals of length $\le2\delta(1+\kappa)v$,
  which enter the difference of the counts with opposite signs; so the counts differ by at most
  $2\delta(1+\kappa)v+1$, and the normalisations $\kappa v$, $\kappa v'$ change the quotient by at most $2.01\delta(1+1/(\kappa v))$,
  which is smaller; and $\rhodet(v)+\rhodet(v')\le\Psi(2v,X_0)/(\kappa v)+\Psi(2v,X_0)/(\kappa v')\le3\Psi(2v,X_0)/(\kappa v)$, as
  $(1+\kappa)v'\le2v$. Put $u:=\log v/\Lambda$
  and $U_0:=\Lambda$. Use the first bound for $u<U_0$ and the second beyond; for $u\ge U_0$, Rankin's bound (exponent
  $1-1/(2\Lambda)$) gives $\Psi(2v,X_0)\le10v\Lambda^{1.65}e^{-u/2}$. With the weights of Corollary~\ref{cor:saias}
  ($\sum_{v<X_0^{U_0}}\log(2v)\le1.5(U_0\Lambda)^2$) and $1/\kappa\le X_0^{5/12}$, $\log x$ times this part of $\rhosh(x)-\rhosh(y)$
  is at most
  \begin{multline*}
    C_S\Big[1.5\Lambda^4(4.0001\,\delta+2/X_0)/\kappa+\sum_{u\ge\Lambda}\log(2v)\cdot30\Lambda^{1.65}e^{-u/2}/\kappa\Big]\\
    \le C_S\big[7\Lambda^4X_0^{-1/12}+90\Lambda^{4.7}e^{-\Lambda/12}\big]\le e^{-10^6},
  \end{multline*}
  using $\delta\le X_0^{-1/2}$. Both terms decrease with $\Lambda$, so this holds for every $\Lambda\ge10^8$; at $\Lambda=10^8$ (with
  $C_S\le77$) the logarithm of the bound is $\le-8.33\cdot10^6$.
\item \emph{Resolvable $R_j'$, smooth change:} by \eqref{eq:Rk} and \eqref{eq:ampl},
  $|R_j''(w)|\le\alpha_j\gamma_j^2w^{\theta-2}(1+O(1/\gamma_j))$, so at a level with $|v-v'|\le2\delta v$,
  $|R_j'(v)-R_j'(v')|\le1.01\,\delta\gamma_j\kappa_j(v/x_j)^{\theta-1}$. Summed with the weights $\mu^2(n)/n$ over $n\le x/x_j$, this is at
  most $1.01F_\theta\,\delta\gamma_j\kappa_j\le2\delta\gamma_j\kappa_j/(1-\theta)\le2\kappa_j/((1-\theta)\gamma_j\log^4x)$, as $F_\theta\le1.33$
  and $\delta\gamma_j^2\le\log^{-4}x$ for resolvable $\gamma_j\le\gamma_*$. This is summable over $j$.
\item \emph{$\sigma_{\Abs}$:} $|\sigma_{\Abs}|\le2L_v^{-4}$, and it vanishes unless $\log\log v>2\log\Lambda+5.97$. So the
  weighted sum is at most $4C_S\sum_v\log(2v)L_v^{-4}/\log x\le C_Se^{-80}/\log x$.
\end{itemize}

% Logged correction 84 (CORRECTIONS-vs-v3.md): v3 defines k := max{j : x_j <= x}, which is undefined for x < x_1, a range
% the proposition covers ([X_sw, x_1) is non-empty); there R' = 0 on [1,x] and the step is void. The case k = 1 (gamma_* = 1) is stated.
\emph{Step 4: the current oscillation, if it is not resolvable.} If $x<x_1$, every $R_j$ vanishes on $[1,x]$ (as $R_j=0$ on
$[1,x_j^+]$), so $R'=0$ there, the oscillation contributes nothing, and this step is void; this happens on the whole range
$[X_{\mathrm{sw}},x_1)$, since $\log X_{\mathrm{sw}}=X_0^{0.01}<X_0^2\le\log x_1$. Otherwise $k\ge1$ is as above; $R_j$ vanishes on
$[1,x]$ for $j>k$. Put $\rho^{(k)}_{\mathrm{osc}}(y):=\sum_nc_n(y)R_k'(y/n)/n$.
\begin{itemize}
\item \emph{Separation.} If $\gamma_k$ is not resolvable, then $\gamma_k>\Gamma(x)$. Every $j<k$ is resolvable (and no $\gamma_j$ with $j>k$ is, as
  $\log\gamma_j\ge\log\Gamma_j\ge(1-\theta)\log x_j-\Lambda-3\log\log x_j-1>c_{\mathrm{LP}}\sqrt{\log x}/4$ for $x_j>x$), and
  $\gamma_*\le\max(1,\gamma_{k-1})\le\log x$ (for $k\ge2$, $\gamma_{k-1}\le x_{k-1}\le\log x$ as $x_k\ge\exp(x_{k-1})$; for $k=1$ there is
  no resolvable $\gamma_j$, and $\gamma_*=1$). So
  $\gamma_k\delta\ge\Gamma(x)\min(\log^{-6}x,X_0^{-1/2})\ge e^{c_{\mathrm{LP}}\sqrt{\log x}/5}$, since, for
  $\sqrt{\log x}\ge X_0^{0.005}=e^{\Lambda/200}$,
  \[
    c_{\mathrm{LP}}\sqrt{\log x}/20\ge(8.8\cdot10^{-4}/\Lambda)\sqrt{\log x}\ge\max(6\log\log x,\Lambda/2).
  \]
% Logged correction 70: the comparison of correction 69 displayed (it overflowed the line).
% Logged correction 69: the comparison made explicit (was "c_LP sqrt(log x) >> log X_0 + log log x").
\item \emph{Average.} Where $c_n(y)$ is constant on $J$, $\int_JR_k'(y/n)\,dy/n=R_k(x(1-\delta)/n)-R_k(x(1-2\delta)/n)$, and
  $|R_k(v)|\le2\kappa_kv(v/x_k)^{\theta-1}/((1-\theta)\gamma_k)$ (integrate the oscillation at frequency $\gamma_k$). So, with
  Lemma~\ref{lem:lemmaM}, the $J$-average is at most $4\kappa_k(\log x+B_1)/((1-\theta)\gamma_k\delta)$, plus the threshold
  part $\le\kappa_k\cdot R(2\delta)$, where $R(\cdot)$ is the bound of Lemma~\ref{lem:threshold-mass}. Both are $o(1/\log x)$.
\item \emph{Pointwise.} $|c_n|\le\mu^2(n)$, and $R_k'(x/n)\ne0$ forces $n\le x/x_k$. So
  \[
    |\rho^{(k)}_{\mathrm{osc}}(x)|\le\kappa_k(x/x_k)^{\theta-1}\sum_{n\le x/x_k}\mu^2(n)n^{-\theta}\le\kappa_kF_\theta.
  \]
  \begin{itemize}
  \item For $x\le x_k^2$: $\kappa_kF_\theta\log x\le4\kappa_k'F_\theta=4\kappa_0F_\theta/\log^2x_k=o(1)$.
  \item For $x_k^2<x<\exp(C'\log^2x_k)$, with $C':=(4(1-\theta)/c_{\mathrm{LP}})^2$:
    \[
      \kappa_kF_\theta\log x\le2\kappa_k'F_\theta C'\log x_k=2\kappa_0F_\theta C'/\log x_k=o(1),
    \]
    since $\log x_k\ge X_0^2$ and $C'\asymp\Lambda^2$.
  \end{itemize}
\item For $x\ge\exp(C'\log^2x_k)$, $\gamma_k\le x_k^{1-\theta}$ is resolvable, and Step~3 covers it.
\end{itemize}

\emph{Step 5.} Combining Steps 1--4,
$\rhosh(x)\log x\ge\pi_0(x)\log x-\eps_N-E_{\mathrm{BR}}-e^{-59}-o(1)\ge1-10^{-12}-o(1)\ge0.97$.
% Logged correction 69 (CORRECTIONS-vs-v3.md): the terms behind o(1) listed with explicit bounds; their sum is
% <= 1.6e-3 Lambda^{-4}, inside the convention o(1) <= Lambda^{-3} of Section sec:noise-notation.
Here $o(1)$ is the sum of the following contributions to $\rhosh(x)\log x$; we use $\log x\ge X_0^{1/100}=e^{\Lambda/100}$,
$1/\delta\le x^{0.02}$, $\delta\le\log^{-4}x$, $\log(1/\delta)\le0.6\sqrt{\log x}$ and $H(x)\le x^{0.01}$.
\begin{itemize}
\item Step~1: $\le x^{-0.4}$ from the partial boxes and the local count; $\le e^{-0.4c_{\mathrm{LP}}\sqrt{\log x}}$ from
  Theorem~\ref{thm:lp}(ii) (as $\log K\le98+5.9\log\Lambda$); and $\le3\delta/\log x$ from the variation of $\pi_0$ over $J$.
\item Step~2: $\le129\log^{-5}x$ from $\bar\nu$ (\S\ref{sec:noise-perscale}); $\le55e^{-0.99\Lambda}$ from $\NT$
  (Lemma~\ref{lem:tau-near}); and $\le x^{-0.9}$ from the boundary term.
% Logged correction 93: the rho_det and R'_j items are at displacement 2 delta; the bounds X_0^{-0.08} and log^{-2} x still hold.
\item Step~3: $\le160\log^{-1.5}x$ from the threshold changes; $\le\log^{-3}x$ from the jump at $X_0$; $\le X_0^{-0.08}$ from
  $\rhodet$ (the bound of Step~3, at displacement $2\delta$, is $\le C_S[7\Lambda^4X_0^{-1/12}+90\Lambda^{4.7}e^{-\Lambda/12}]\le X_0^{-0.08}$);
  $\le\log^{-2}x$ from the resolvable $R_j'$ (namely $\le4\sum_j\kappa_j\log^{-3}x$, as $1-\theta\ge\frac12$); and $\le4C_S\sum_{L_v>391\Lambda^2}\log(2v)L_v^{-4}\le1.5\cdot10^{-3}\Lambda^{-4}$
  from $\sigma_{\Abs}$.
\item Step~4: $\le e^{-c_{\mathrm{LP}}\sqrt{\log x}/6}+40\log^{-1.5}x$ from the average, and $\le2\cdot10^4\Lambda^2X_0^{-2}$ from the
  pointwise bounds (with $\kappa_0\le\frac14$, $F_\theta\le1.33$ and $C'\le(171\Lambda)^2$).
\end{itemize}
In total $o(1)\le1.5\cdot10^{-3}\Lambda^{-4}+e^{-0.01\Lambda}\le1.6\cdot10^{-3}\Lambda^{-4}$, for every $\Lambda\ge10^8$.
\end{proof}

\begin{remark}
No twisted M\"obius estimate is used: old oscillations are \emph{resolved} by the window of the local prime number theorem,
not cancelled. Range (A) uses no local prime number theorem, and range (B) uses no Huxley-type input beyond
$\jz\le1.0001/\Lambda$. The ranges fit: $\log X_{\mathrm{LP}}=10^5\Lambda^2\log^2\Lambda\le\log X_{\mathrm{sw}}=X_0^{1/100}\le X_0^{0.079}
\le(0.0048/\kappa)^{1/5}$.
% Logged correction 69: the ranges made explicit (were "<<").
\end{remark}
